\section{Spectra of $\Ctwo$ formulas} \label{sec:ctwo}

%In this section, our goal will be to show:

%\%begin{theorem} \label{thm:ctwoandglobal}
%    The finite satisfiability problem of $\Ctwo$ with global constraints is in $\nexp$.
%\end{theorem}

In this section we will consider converting formulas in logic with arithmetic into linear constraints. These constraints will characterize  the set of models of a $\Ctwo$ formula. With this characterization in place, we have reduced finite satisfiability of the formula to a check for satisfiability of the conjunction of the linear constraints, which can be done via standard linear algebra solving. From this it will follow that $\Ctwo$ with a variety of global cardinality constraints is decidable: we just conjoined the cardinality constraints with the constraints characterizing the models.

We fix a vocabulary $\vocab$ which consists of $n$ unary predicates $U_1, \ldots, U_n$ and $m$ binary predicates $R_1, \ldots, R_m$.
We consider $\Ctwo$ sentence $\varphi$ over $\vocab$ in \emph{normal form}:
\begin{equation*}
    \varphi\ :=\ 
    \forall x\ \gamma(x)\ \land\ 
    \forall x\ \forall y\ \left(x \neq y \to \alpha(x, y)\right)\ \land\ 
    \bigwedge_{i \in \intsinterval{m'}} \forall x\ \exists^{=k_i} y\ \left(R_i(x,y) \land x \neq y\right),
\end{equation*}
where $\gamma(x)$ and $\alpha(x, y)$ are quantifier-free,
$R_i$ are binary predicates,
$0 \le m' \le m$,
and $k_i \in \bbN$.

\begin{definition}[Count bound vector]
The count bound vector of $\varphi$, denoted by $\bfk^\varphi$, 
is the $m'$ dimension vector such that
the $i^{th}$ component of $\bfk^\varphi$ is $k_i$.
\end{definition}

We denote the sum of $k_i$ by $K^\varphi$.
\begin{definition}[Forward and backward characteristic vector of 2-types]
    For $m'$ as above and a 2-type $\eta$, the forward characteristic vector of $\eta$,
    denoted by $\eta^\trif$, is an $m'$ dimensional vector defined:
    for every $1 \le i \le m'$, if $R_i(x, y) \in \eta$,
    then the $i^{th}$ component of $\eta^\trif$ is 1.
    Otherwise, it is $0$.
    The backward characteristic vector of $\eta$,
    denoted by $\eta^\trib$,
    is defined similarly.
\end{definition}

\begin{definition}[Forward and backward silent 2-type]
    For $m'$ as above and a 2-type $\eta$ 
    we say that $\eta$ is forward-silent if for $1 \le i \le m'$,
    $\neg R_i(x, y) \in \eta$.
    Backward-silent is defined similarly.
\end{definition}

We define the notions of $\Onetypes^\varphi$ and $\twotypes^\varphi_{\pi_1, \pi_2}$ similarly to those in $\GPtwo$.
We also define
$\cKvpnull_{\pi_1, \pi_2} \subseteq \twotypes^\varphi_{\pi_1, \pi_2}$ as the set of 2-types that are both forward- and backward-silent,
$\cKvpf_{\pi_1, \pi_2}$ for 2-types that are not forward-silent but are backward-silent,
and
$\cKvpfb_{\pi_1, \pi_2}$ for 2-types that are neither forward- nor backward-silent.

In the rest of the section, we deal only with normal form sentences. We explain how the results apply to arbitrary $\Ctwo$ sentences in the appendix.

\begin{definition}[Spectrum and $1$-type spectrum]
    For every $\Ctwo$ sentence $\varphi$,
    the spectrum of $\varphi$, denoted by $\spec(\varphi)$,
    is the set of natural numbers that are sizes of finite models of $\varphi$.

    The 1-type spectrum of $\varphi$, denoted by $\pispec(\varphi)$,
    is the set of 1-type cardinality vectors (see Definition~\ref{def:cardboundvector}) corresponding to finite models of $\varphi$.
\end{definition}

We aim to show: 

\begin{theorem} \label{thm:spectrasemilinear}
    For every $\Ctwo$ sentence $\varphi$,
    $\spec(\varphi)$ and $\pispec(\varphi)$ are semilinear.
    Furthermore, each spectrum can be represented by an existential Presburger formula of size doubly exponential in $\varphi$.
\end{theorem}

\begin{definition}[Silent compatible 1-types]
    For a pair of 1-types $\pi_1$ and $\pi_2$,
    we say that $\pi_1$ and $\pi_2$ are \emph{silent-compatible} w.r.t. $\varphi$,
    if $\cKvpnull_{\pi_1, \pi_2}$ is non-empty. Otherwise, we say, following~\cite{ctwocomplexity}, that 
    $\pi_1$ and $\pi_2$ are a \emph{noisy pair}.
\end{definition}

\begin{lemma} \label{lem:differentiated}
    For every $\Ctwo$ sentence $\varphi$ in normal form and $\vocab$-structure $\cG \models \varphi$,
    for every pair of 1-types $\pi_1$ and $\pi_2$,
    if $\pi_1$ and $\pi_2$ are a noisy pair (w.r.t. $\varphi$), then
    at least one of the 1-types $\pi_1$ or $\pi_2$ is realized by at most $2K^\varphi+1$ vertices in $\cG$.
\end{lemma}

\begin{proof}
    
    Recall that $V_{\pi}$ denotes the vertices in $V$ with 1-type $\pi$.
    Because $\cG \models \varphi$, for vertex $v \in \cG$,
    $\sum_{u \in V \setminus \set{v}} \twotypeof^\trif(v, u) = \bfk^\varphi$.
    We consider two cases.

    \textbf{\uline{Suppose that $\pi_1 = \pi_2$.}}
    Since $\twotypeof^\trif(v, u) = \twotypeof^\trib(u, v)$,
    we have
    \begin{equation*}
        \begin{aligned}
            \sum_{v \in V_{\pi_1}} \sum_{u \in V_{\pi_1} \setminus \set{v}}
            (\twotypeof^\trif(v, u) + \twotypeof^\trib(v, u))
            \ =\ &
            2 \sum_{v \in V_{\pi_1}} \sum_{u \in V_{\pi_1} \setminus \set{v}}
            \twotypeof^\trif(v, u) \\
            \ \vle\ &
            2 \sum_{v \in V_{\pi_1}} \sum_{u \in V \setminus \set{v}} \twotypeof^\trif(v, u)
            \ =\ 2\abs{V_{\pi_1}} \bfk^\varphi.
        \end{aligned}
    \end{equation*}
    For every pair of vertices $v$ and $u$ in $V_{\pi_1}$ with $v \neq u$,
    by definition of noisy pair,
    the sum of components of $\twotypeof^\trif(v, u) + \twotypeof^\trib(v, u)$ is at least $1$.
    Therefore, we have
    \begin{equation*}
        \sum_{v \in V_{\pi_1}} \sum_{u \in V_{\pi_1} \setminus \set{v}} 1
        \ =\ \abs{V_{\pi_1}}\cdot(\abs{V_{\pi_1}} - 1)
        \ \le\ 2\abs{V_{\pi_1}} \cdot K^\varphi,
    \end{equation*}
    which implies that $\abs{V_{\pi_1}} \le 2K^\varphi + 1$. %\michael{$2K$ should be $2K^\varphi$?}\ch{yes. thanks}
    
    \textbf{\uline{Suppose that $\pi_1 \neq \pi_2$.}}
    By a similar argument, 
    $\abs{V_{\pi_1}} \cdot \abs{V_{\pi_2}} \le (\abs{V_{\pi_1}} + \abs{V_{\pi_2}}) \cdot K^\varphi$.
    If $\abs{V_{\pi_1}} \leq 2K^\varphi+1$ we are done. So assume the opposite.
    Then we have
    \begin{equation*}
        \abs{V_{\pi_2}}
        \ \le\ 
        \frac{K^\varphi}{1-\frac{K^\varphi}{\abs{V_{\pi_1}}}}
        \ \le\ 
        \frac{K^\varphi}{1-\frac{K^\varphi}{2K^\varphi + 1}}
        \ \le\ 
        2K^\varphi.
    \end{equation*}
\end{proof}

The following straightforward proposition gives
conditions that characterize the satisfiability of $\Ctwo$ by a model in terms of the prior constructs.
\begin{proposition}\label{prop:ctwo_graph}
    For every $\Ctwo$ sentence $\varphi$ in normal form and $\vocab$-structure $\cG$,
    $\cG$ is a model of $\varphi$ if and only if following conditions hold.
    For every $v, u \in V$ with $v \neq u$,
    \begin{compactitem}
        \item (1-type) $\onetypeof(v) \in \Onetypes^\varphi$,
        \item (2-type) $\twotypeof(v, u) \in \twotypes^\varphi_{\onetypeof(v), \onetypeof(u)}$, and
        \item (Counting) $\sum_{u \in V \setminus \set{v}} \twotypeof^\trif(v, u) = \bfk^\varphi$,
    \end{compactitem}
    where $\bfk^\varphi$ is the count bound vector of $\varphi$.
\end{proposition}

\begin{definition}[Partial model of $\Ctwo$ sentence] \label{def:partial}
    For every $\Ctwo$ sentence $\varphi$ and $\vocab$-structure $\cG$,
    we say that a complete colored graph $\cG$ is a \emph{partial model} of $\varphi$
    if $\cG$ satisfies the first two conditions in Proposition~\ref{prop:ctwo_graph}.
\end{definition}

The following is obvious.
\begin{lemma}\label{lem:core_partial_model}
    For every $\Ctwo$ sentence $\varphi$ in normal form,
    for every $\vocab$-structure $\cG \models \varphi$,
    for every $\cH \subseteq \cG$,
    $\cH$ is a partial model of $\varphi$.
\end{lemma}

\myparagraph{Core of structures}
\begin{definition} \label{def:core}
    For $\cG$ a $\vocab$-structure, $\ell \in \bbN$,
    we say that substructure $\cH \subseteq \cG$ is an $\ell$-core of $\cG$ if 
    \begin{compactitem}
        \item for every 1-type $\pi$,
        the number of vertices in $\cG \setminus \cH$
        with 1-type $\pi$ is either $0$ or at least $\ell$
        \item for  1-types $\pi_1$, $\pi_2$, and 2-type $\eta$,
        there are either $0$ and at least $\ell$
        pairs of vertices $v, u \in \cG \setminus \cH$ satisfying that
        $v \neq u$,
        $\onetypeof(v) = \pi_1$,
        $\twotypeof(v, u) = \eta$,
        and $\onetypeof(u) = \pi_2$.
    \end{compactitem}
\end{definition}

Note that the core of $\cG$ is not unique,
and $\cG$ is trivially a core of itself for every $\ell \in \bbN$.
Therefore, we are interested in the existence of a \emph{small} core of $\cG$.
We can show that for every $\ell$, every $\vocab$-structure has an \emph{exponential-size} core.

\begin{lemma}\label{lem:small_core}
    For every $\vocab$-structure $\cG$, for every $\ell \in \bbN$,
    $\cG$ has an $\ell$-core $\cH$ with size at most $(2^{2n+4m+2}) \cdot \ell$.
\end{lemma}

The intuition behind the construction is to repeatedly select a 1-type or tuple that violates the $\ell$-core condition
and add the relevant vertices to $\cH$.
Since there are only finitely many such choices, the process eventually stabilizes,
and the size of the constructed core remains bounded: see the appendix for details.
 
For a $\Ctwo$ sentence $\varphi$ and a $\vocab$-structure $\cG \models \varphi$,
if $\cH$ is a $\ell$-core of $\cG$ with sufficiently large $\ell$,
then we have the following important property of $\cG \setminus \cH$:

\begin{lemma}\label{lem:mutually_null_compatible}
    For every $\Ctwo$ sentence $\varphi$ and $\vocab$-structure $\cG \models \varphi$,
    for every $\ell$-core $\cH \subseteq \cG$ with $\ell > 2K^\varphi + 1$, 
    the 1-types realized in $\cG \setminus \cH$ are silent-compatible (w.r.t. $\varphi$).
\end{lemma}

\begin{proof}
    Suppose that there exist 1-types $\pi_1$ and $\pi_2$ realized in $\cG \setminus \cH$
    which are not  silent-compatible.
    By Lemma~\ref{lem:differentiated},
    at least one of the 1-types $\pi_1$ or $\pi_2$ is realized by at most $2K^\varphi+1$ vertices in $\cG$.
    Without loss of generality,
    suppose that $\pi_1$ is realized by at most $2K^\varphi+1$ vertices in $\cG$,
    then there are at most $2K^\varphi+1 < \ell$ vertices in $\cG \setminus \cH$ realized $\pi_1$.
    However, by the definition of $\ell$-core, since $\pi_1$ is realized in $\cG \setminus \cH$,
    it is realized by at least $\ell$ vertices in $\cG \setminus \cH$.
    This leads to a contradiction.
\end{proof}

In Section~\ref{sec:gptwo}, we introduced the notion of behavior vectors to describe the number of edges connected to a given vertex with specific 2-types.
In this section, we consider $\vocab$-structures with a fixed core and extend the concept of behavior vectors to incorporate information about the core.
Again, we fix an ordering of the set $\twotypes^+ \times \Onetypes$,
which has cardinality $2^{n+3m}-2^{n+m}$.

\begin{definition}[Extended behavior]
    Fixing a vocabulary $\vocab$, an \emph{extended behavior} w.r.t. 
    a set of vertices $V$ is a tuple $\tuple{g, \bff}$ where
    $g$ is a mapping from a vertex $v \in V$ to a 2-type in $\twotypes$, 
    and
    $\bff$ is an element of $\bbN^{2^{n+3m}-2^{n+m}}$.

\end{definition}

Informally, each extended behavior consists of two components.
The first one is a mapping from vertices in the core to 2-types,
capturing the 2-type realized by the edge between a given vertex and core vertices.
The second component is a behavior vector.

Abusing notation, for an extended behavior $\tuple{g, \bff}$,
we will write $g^\trif(v)$ for ${(g(v))}^\trif$ and
$\dual{g}(v)$ for the dual of $g(v)$.
We will also write
$\bff(\eta, \pi)$ for the $i^{th}$ component of $\bff$,
where $i$ is the index of the tuple $\tuple{\eta, \pi}$ in the fixed ordering of $\twotypes^+ \times \Onetypes$.

\begin{definition}[Compatibility of extended behaviors]\label{def:compatextended}
    For a $\Ctwo$ sentence $\varphi$ in normal form,
    $\pi \in \onetypes^\varphi$, and
    a $\vocab$-structure $\cH$ with vertices $V$,
    we say that an extended behavior function $\tuple{g, \bff}$ w.r.t. $V$ is
    \emph{compatible} with $\varphi$, $\pi$, and $\cH$
    if $f$ satisfies:
    \begin{enumerate}
        \item (Compatibility with $\cH$)
        For every $v \in V$, $g(v) \in \twotypes^\varphi_{\pi, \onetypeof(v)}$.

        \item (Compatibility with 1- and 2- types)
        For every $\tuple{\eta, \pi'} \in \twotypes^+ \times \Onetypes$,
        if $\pi' \not\in \onetypes^\varphi$ or 
        $\eta \not\in \twotypes^\varphi_{\pi, \pi'}$,
        then $\bff(\eta, \pi') = 0$.

        \item (Forward-silent zero)
        For every $\tuple{\eta, \pi'} \in \twotypes^+ \times \Onetypes$,
        if $\eta$ is forward-silent,
        then $\bff(\eta, \pi') = 0$.
       
        \item (Counting) The remaining entries of $\bff$,
        that is, entries for $2$-types that are not forward-silent, are bounded by the following constraint:
        \begin{equation*}
            \sum_{v \in V} g^\trif(v) + 
            \sum_{\pi' \in \Onetypes^\varphi}
            \left(
                \sum_{\eta \in \cKvpf_{\pi, \pi'}} \bff(\eta, \pi') \cdot \eta^\trif +
                \sum_{\eta \in \cKvpfb_{\pi, \pi'}} \bff(\eta, \pi') \cdot \eta^\trif
            \right)
            \ =\ \bfk^\varphi.
        \end{equation*}
        Note that this condition implies that
        if $\pi'$ and $\eta$ are compatible with $\varphi$, and
        $\eta$ is not forward-silent,
        then $\bff(\eta, \pi')$ is bounded by $K^\varphi$.
    \end{enumerate}
\end{definition}

Note that given an extended behavior,
we can check whether it is compatible with $\pi$ and $\cH$ in time polynomial in $2^n$ and $\abs{V}$.
Let $\cB^{\varphi}_{\pi, \cH}$ be the set of extended behaviors w.r.t. 
the vertices of $\cH$ that are compatible with $\varphi$, $\pi$ and $\cH$.
For each $\tuple{g, \bff} \in \cB^{\varphi}_{\pi, \cH}$,
recall that every entry of $\bff$ is bounded by $\abs{K^\varphi}$.
Thus the cardinality of $\cB^{\varphi}_{\pi, \cH}$ is bounded by
$\abs{V}^\abs{\cK} \cdot \abs{K^\varphi}^\abs{\twotypes \times \onetypes}$
which is doubly exponential in the length of $\varphi$.

\myparagraph{Computing a semilinear representation of the spectrum of a $\Ctwo$ sentence}
We are now ready to show that the $1$-type spectrum of a $\Ctwo$ sentence $\varphi$ in normal form is semilinear.
For every finite model $\cG$ of $\varphi$,
by Lemma~\ref{lem:differentiated},
for sufficiently large $\ell$,
for every $\ell$-core $\cH$ of $\cG$
and pair of 1-types $\pi_1$ and $\pi_2$ in $\cG \setminus \cH$,
$\pi_1$ and $\pi_2$ are silent-compatible w.r.t. $\varphi$. 
Furthermore by Lemma~\ref{lem:small_core},
$\cG$ has an $\ell$-core whose size is exponential in $\varphi$ and $\ell$.
We prove Theorem~\ref{thm:c2_spec} via a two-step construction.
First we show that for every such core $\cH$,
there exists an existential Presburger formula $\Psi^\varphi_\cH(\bfx)$
such that the solution set of $\Psi^\varphi_\cH(\bfx)$ corresponds to the set of 1-type cardinality vectors of finite models of $\varphi$ with $\ell$-core $\cH$.

\begin{compactitem}
    \item The variables are $x_\pi$ for $\pi \in \onetypes^\varphi$,
    and $y_{\pi, \tuple{g, \bff}}$ for $\pi \in \onetypes^\varphi$ and $\tuple{g, \bff} \in \cB^\varphi_{\pi, \cH}$.
    Intuitively, $\bfx$ is the 1-type spectrum of the finite model $\cG$ and $y_{\pi, \tuple{g, \bff}}$ is the number of vertices in $\cG \setminus \cH$ with 1-type $\pi$ and extended behavior $\tuple{g, \bff}$.
    \item
    The formula $\summ_{\cH}(\bfx, \bfy)$ asserts that the 1-type spectrum of $\cG$ is the sum of the 1-type spectrum of $\cH$ and
    the $1$-type spectrum of $\cG \setminus \cH$.
    Let $s_\pi$ be the number of vertices in $\cH$ with 1-type $\pi$. \\
    \begin{notsotiny}
    $\displaystyle\summ_{\cH}(\bfx, \bfy)
         := 
     \bigwedge_{\pi \in \onetypes^\varphi} \left(x_\pi = s_\pi + \sum_{\tuple{g, \bff} \in \cB^{\varphi}_{\pi, \cH}} y_{\pi, \tuple{g, \bff}}\right)$.
    \end{notsotiny}
    \item
    The formula $\comp_{\cH}(\bfy)$ guarantees that the vertices in $\cH$ satisfies the Counting condition of Proposition~\ref{prop:ctwo_graph}.
    Let $V_c$ be the vertices in $\cH$. \\
    \begin{notsotiny}
    $\displaystyle\comp_{\cH}(\bfy)
         := 
        \bigwedge_{v \in V_c} \left(
            \sum_{u \in V_c \setminus \set{v}} \twotypeof^\trif(v, u) +
            \sum_{\pi \in \onetypes^\varphi} \sum_{\tuple{g, \bff} \in \cB^{\varphi}_{\pi, \cH}} g^\trib(v) \cdot y_{\pi, \tuple{g, \bff}}
        = \bfk^\varphi \right)$.
        \end{notsotiny}
    \item
    The formula $\silent(\bfy)$ assert that 1-types realized in $\cG \setminus \cH$ are silent compatible. \\
    \begin{notsotiny}
    $\displaystyle\silent_{\cH}(\bfy)
         := 
        \bigwedge_{\substack{\pi_1, \pi_2 \in \onetypes^\varphi \\ \text{are not silent compatible}}}
        \left(\sum_{\tuple{g, \bff} \in \cB^{\varphi}_{\pi_1, \cH}} y_{\pi_1, \tuple{g, \bff}} = 0\right) \lor
        \left(\sum_{\tuple{g, \bff} \in \cB^{\varphi}_{\pi_2, \cH}} y_{\pi_2, \tuple{g, \bff}} = 0\right)$.
        \end{notsotiny}
    \item
    The formula $\bige_{\cH}(\bfy)$ guarantees that 1-types and tuples realized in $\cG \setminus \cH$ satisfy the requirements of a core. \\
    \begin{notsotiny}
    $
    \begin{aligned}
            \bige_{\cH}(\bfy)
        \ :=\ &
        \bigwedge_{\pi \in \onetypes^\varphi}
            \exists z\ 
            \left(
                z = \sum_{\tuple{g, \bff} \in \cB^\varphi_{\pi, \cH}} y_{\pi, \tuple{g, \bff}}
            \right) \land
            \left(
                z = 0  \lor
                z \ge {(2K^\varphi+1)}^2 
            \right)
        \land \\
        &\bigwedge_{\substack{
            \pi_1, \pi_2 \in \onetypes^\varphi \\ 
            \eta \in \cKvpfb_{\pi_1, \pi_2}}}
        \exists z\ 
        \left(z = \sum_{\tuple{g, \bff} \in \cB^\varphi_{\pi_1, \cH}} \bff(\eta, \pi_2) \cdot y_{\pi_1, \tuple{g, \bff}}\right) \land
        \left(
            z = 0 \lor
            z \ge (2K^\varphi+1)^2
            \right).
        \end{aligned}
    $
    \end{notsotiny}
    \item
        The formula $\match_{1, \cH}(\bfy)$ encodes the edge matching condition for audible 2-types and vertices in $\cG \setminus \cH$. \\
        \begin{notsotiny}
        $\displaystyle
        \match_{1, \cH}(\bfy)
         := 
        \bigwedge_{\substack{
            \pi_1, \pi_2 \in \onetypes^\varphi \\
            \eta \in \cKvpfb_{\pi_1, \pi_2}}}
        \left(
            \sum_{\tuple{g, \bff} \in \cB^{\varphi}_{\pi_1, \cH}} \bff(\eta, \pi_2) \cdot y_{\pi_1, \tuple{g, \bff}}
            =
            \sum_{\tuple{g, \bff} \in \cB^{\varphi}_{\pi_2, \cH}} \bff(\dual{\eta}, \pi_1) \cdot y_{\pi_2, \tuple{g, \bff}}
        \right)
        $.
        \end{notsotiny}

    \item
        The formula $\match_{2, \cH}(\bfy)$ encodes the edge matching condition for backward-silent but not forward-silent 2-types and vertices in $\cG \setminus \cH$. \\
        \begin{notsotiny}
        $\displaystyle
        \match_{2, \cH}(\bfy)
         := 
        \bigwedge_{\substack{
            \pi_1, \pi_2 \in \onetypes^\varphi \\
            \eta \in \cKvpf_{\pi_1, \pi_2}}}
        \left(
            \left(
                \sum_{\tuple{g, \bff} \in \cB^{\varphi}_{\pi_1, \cH}} \bff(\eta, \pi_2) \cdot y_{\pi_1, \tuple{g, \bff}} > 0
            \right)
            \to
            \left(
                \sum_{\tuple{g, \bff} \in \cB^{\varphi}_{\pi_2, \cH}} y_{\pi_2, \tuple{g, \bff}} > 0
            \right)
        \right)
        $.
        \end{notsotiny}
        Keep in mind that we will always be considering solutions in the natural numbers. So this implication could be rewritten without summation: if for one of the extended behaviors for this triple, $\bff(\eta, \pi_2) \cdot y_{\pi_1, \tuple{g, \bff}} > 0$, then  one of the numbers $y_{\pi_2, \tuple{g, \bff}} > 0$.
\end{compactitem}
Finally, $\displaystyle
    \Psi^\varphi_{\cH}(\bfx)
         :=
         \exists \bfy\ 
        \summ_{\cH}(\bfx, \bfy) \land
        \comp_{\cH}(\bfy) \land
        \silent_{\cH}(\bfy) \land
        \bige_{\cH}(\bfy) \land
        \bigwedge_{i=1, 2}
        \match_{i, \cH}(\bfy)$.

This system characterizes the spectrum:
\begin{lemma}\label{lem:c2_spec_core}
    For every $\Ctwo$ sentence $\varphi$ and partial model $\cH$ of $\varphi$,
    a vector $\bfv$ is a solution of $\Psi^\varphi_\cH(\bfx)$
    if and only if
    there exists a finite model $\cG$ of $\varphi$ such that
    the 1-type cardinality vector of $\cG$ is $\bfv$,
    and $\cH$ is a $(2K^\varphi+1)^2$-core of $\cG$.
\end{lemma}

It is easy to see that vectors in the spectrum satisfy the equations. 
In the other direction, we extend the core $\cH$ to a finite model $\cG$ of $\varphi$.
Solutions to the equations tell us how many vertices of each $1$-type and extended behavior to create in $\cG$.
We assign edge colors so that each vertex realizes its intended extended behavior.
This is always possible because the solution is sufficiently large, as ensured by the Big condition in the equation. This extra size gives us some leeway to match up $1$-types via $2$-types.
Details are in the appendix.

Let $\mathfrak{H}$ be the collection of partial models of $\varphi$ with size at most $2^{2n+4m+2} \cdot (2K^\varphi+1)^2$.
We define the existential Presburger formula 
$\Psi^\varphi(\bfx)
 := 
\bigvee_{\cH \in \mathfrak{H}}
\Psi^\varphi_{\cH}(\bfy).$

\begin{lemma}\label{thm:c2_spec}
    For every $\Ctwo$ sentence $\varphi$,
    $\sol{\Psi^\varphi(\bfx)} = \pispec(\varphi)$.
\end{lemma}

\begin{proof}
    Suppose that $\bfv \in \sol{\Psi^\varphi(\bfx)}$,
    there exists $\cH$ such that $\bfv \in \sol{\Psi^\varphi_\cH(\bfx)}$.
    By Lemma~\ref{lem:c2_spec_core},
    there exists a model $\cG$ of $\varphi$ such that the 1-type cardinality vector of $\cG$ is $\bfv$.
    Thus $\bfv \in \pispec(\varphi)$.
    On the other hand, suppose that $\bfv \in \pispec(\varphi)$.
    Then there is a model $\cG$ of $\varphi$ with 1-type cardinality vector  $\bfv$.
    By Lemma~\ref{lem:small_core}, $\cG$ has a $(2K^\varphi+1)^2$-core $\cH$ with size at most $2^{2n+4m+2}(2K^\varphi+1)^2$,
    which implies that $\cH \in \mathfrak{H}$.
    By Lemma~\ref{lem:c2_spec_core},
    $\bfv \in \sol{\Psi^\varphi_\cH(\bfx)}$.
    Since $\cH \in \mathfrak{H}$,
    $\bfv \in \sol{\Psi^\varphi(\bfx)}$.
\end{proof}

Theorem~\ref{thm:spectrasemilinear} follows immediately from Lemma~\ref{thm:c2_spec}.

%\ch{added:}
Note that the length of the formula $\Psi^\varphi(\bfx)$ is doubly exponential in the length of $\varphi$,
since the length of $\Psi^\varphi_\cH(\bfx)$ is already doubly exponential in the length of $\varphi$, 
as the number of variables $y_{\pi, \tuple{g, \bff}}$ is doubly exponential.
Moreover, there are doubly exponentially many possible cores $\cH$.
%\michael{Broken sentence?}
%\ch{the following is our TODO, I commented out it.}
% In order to obtain an exponential length Presburger formula that characterizes the spectrum, we encode the core $\cH$ by exponentially many Boolean variables.
% Let $\bfz^{node}$ be the tuple of variables
% $z^{node}_{i, \pi}$ for $1 \le i \le (2K+1)^2$ and $\pi \in \onetypes$,
% $z^{edge}_{i, j, \eta}$ for $1 \le i, j \le (2K+1)^2$ and $\eta \in \twotypes$.

\myparagraph{Consequences for finite satisfiability of $\Ctwo$ with global constraints}
Note that in Lemma~\ref{thm:c2_spec},
we show that for every $\Ctwo$ sentence $\varphi$ in normal form,
there exists an existential Presburger formula $\Psi^\varphi(\bfx)$ of doubly exponential size 
such that $\pispec{\varphi} = \sol{\Psi^\varphi(\bfx)}$.
Since for every 1-type $\pi$ and partial model $\cH$ of $\varphi$, we can compute $\cB^{\varphi}_{\pi, \cH}$ in doubly exponential time in the length of $\varphi$ and $\cH$ by a straightforward enumerate-and-check procedure,
$\Psi^\varphi(\bfx)$ can be computed from $\varphi$ in doubly exponential time.
Note that global constraints of $\Ctwo$ sentences are linear constraints over the $1$-type cardinality vector.
Therefore the finite satisfiability problem of $\Ctwo$ with global constraints reduces
to satisfiability problem of $\Psi^\varphi(\bfx)$ together with those linear constraints.
This naive approach yields a $2\text{-}\nexp$ algorithm.

We observe that $\Psi^\varphi(\bfx)$ can be rewritten as disjunctions of doubly exponentially many integer linear systems.
Each system contains doubly exponentially many variables, but only \emph{exponentially many constraints}.
By Lemma~\ref{lem:ilp}, if such a system has a solution,
then there exists a solution in which only exponentially many variables are assigned non-zero values.

Therefore, we can decide the finite satisfiability of $\Ctwo$ with global constraints 
using the $\nexp$ procedure in Algorithm~\ref{alg:ctwosat}.

\begin{algorithm}
\begin{algorithmic}[1]
\caption{Algorithm for the finite satisfiability problem of $\Ctwo$ with global constraints}\label{alg:ctwosat}
\Procedure{SAT}{$\varphi$}
    \State Compute $\Onetypes^\varphi$ and $\twotypes^\varphi_{\pi_1, \pi_2}$
    %\State Guess a mutually silent-compatible subset $\Pi' \subseteq \Onetypes^\varphi$
    \State Guess a exponential size partial model $\cH$ of $\varphi$
    \State Guess exponentially many non-zero variables $y'_{\pi, \tuple{g, \bff}}$
    \State Construct the system $\Psi^\varphi_{\cH}(\bfx)$ with only variables $y'_{\pi, \tuple{g, \bff}}$
    \State Check the feasibility of $\Psi^\varphi_{\cH}(\bfx)$ with global constraints
\EndProcedure
\end{algorithmic}
\end{algorithm}

We emphasize again that the argument here is completely generic, and the same comment holds to show that for \emph{any} decidability extension of Presburger arithmetic (e.g. B\"uchi Arithemetic or Semenov Arithmetic~\cite{semenovpresexp}), $\Ctwo$ with global cardinality constraints of that form is decidable, with the complexity the maximum of $\nexptime$ and the complexity of solving the corresponding constraints.

